
Trustworthiness dimensions in large language models are not independent but coupled through training mechanisms. Synchronized evaluation on 817 TruthfulQA prompts with Llama-2-7b demonstrates correlations ranging from r=-0.25 (fairness-reliability) to r=0.72 (reliability-robustness on factual content). Two coupling patterns are validated with mechanistic interpretations: positive reliability-robustness correlation consistent with shared memorization, and negative fairness-reliability correlation consistent with RLHF trade-offs.

The methodological contribution is synchronized evaluation as a paradigm for discovering cross-dimensional correlation structure in existing benchmarks. By treating evaluation logs as multi-dimensional datasets, the work recovers correlation information that was present but discarded in dimension-isolated approaches.

Future work can extend this in three directions: (1) Longitudinal training checkpoint analysis measuring correlation structure across checkpoints (pre-training $\rightarrow$ instruction tuning $\rightarrow$ RLHF) to causally attribute coupling emergence to specific training stages. (2) Cross-architectural generalization replicating h-m1 and h-m2 on GPT-4, Claude, and Gemini to test whether coupling patterns are universal or architecture-specific. (3) Scaling the moderation hypothesis (h-m3) with $n\geq 100$ per stratum to conclusively test whether factual prompts show significantly stronger coupling than misinformation prompts.

The two validated coupling patterns (r=0.72 memorization, r=-0.25 alignment tax) enable actionable insights: practitioners can diagnose memorization strength via reliability-robustness correlation, estimate alignment tax before deploying RLHF, and predict how improving one dimension affects others. This shifts safety intervention design from trial-and-error post-hoc measurement to principled cost-benefit analysis with quantitative estimates.
